\chapter{ML-модели для операционного и финансового управления}
\label{ch:ml}

\section{Архитектура прогнозного управления}

Финансовый аудит, описанный в предыдущей главе, представляет собой ретроспективный анализ: он объясняет, что произошло за прошедшие 12 месяцев. Следующий шаг --- перейти от объяснения прошлого к прогнозированию будущего. Именно этот переход превращает разовый аудит в регулярный управленческий инструмент. Согласно результатам соревнований по прогнозированию M4 и M5~\cite{makridakis2018m4,makridakis2022m5}, на временных рядах малой длины простые методы конкурентоспособны со сложными ML-моделями, что делает выбор метода ключевым методологическим решением.

В рамках проекта разработаны четыре взаимосвязанные ML-задачи, образующие замкнутый контур управления: \textbf{прогноз спроса по SKU} (оперативное планирование производства на 2--4 недели), \textbf{прогноз возвратов} (расчёт производственного плана с учётом возвратного потока), \textbf{прогноз денежного потока} (прогнозирование входящих и исходящих платежей на горизонте 30 дней), \textbf{детектор аномалий} (флаговая маркировка нетипичных операций для ручного контроля). Единой метрикой качества для задач прогнозирования служит WAPE~\cite{hyndman2006another} --- взвешенная средняя абсолютная процентная ошибка, устойчивая к нулевым значениям и пропорционально взвешивающая ошибки по объёму:

\begin{equation}
\text{WAPE} = \frac{\sum_{t=1}^{T} |y_t - \hat{y}_t|}{\sum_{t=1}^{T} y_t}
\end{equation}

Значение WAPE = 1.0 соответствует ошибке, равной средним продажам (прогноз нуля на весь период). WAPE предпочтительнее MAPE при работе с малым числом товаров, поскольку не штрафует бесконечно за нулевые наблюдения~\cite{makridakis2022m5}.

\section{Прогнозирование спроса по SKU}

Данные о продажах получены из отчётов 1С о реализации по контрагентам; временной ряд построен на уровне «SKU~$\times$~месяц» по отгруженным единицам продукции. Активных SKU (продажи $>0$ в $\ge 3$ последних месяцах) --- \textbf{3}: Ажурная коса, Кекс с изюмом, Песочное печенье. Общая валовая выручка по данным 1С: \textbf{53\,270\,000 руб.}, возвраты: \textbf{8\,300\,000 руб.} (15,6\%), обучающая выборка: \textbf{106 строк} (SKU~$\times$~месяц, май 2025 --- март 2026), тестовая: \textbf{24 строки} (апрель 2026). Детальные данные о продажах и возвратах по товарам представлены в таблице~\ref{tab:sales_by_sku}.

\begin{table}[H]
\centering
\caption{Продажи и возвраты по основным товарам}
\label{tab:sales_by_sku}
\begin{tabular}{lrrrr}
\toprule
\textbf{SKU} & \textbf{Валово, руб.} & \textbf{Возвраты, руб.} & \textbf{Нетто, руб.} & \textbf{Return rate} \\
\midrule
Ажурная коса & 18\,200\,000 & 2\,650\,000 & 15\,550\,000 & 14,6\% \\
Кекс с изюмом & 21\,500\,000 & 4\,940\,000 & 16\,560\,000 & 23,0\% \\
Песочное печенье & 13\,570\,000 & 710\,000 & 12\,860\,000 & 5,2\% \\
\midrule
\textbf{Итого} & \textbf{53\,270\,000} & \textbf{8\,300\,000} & \textbf{44\,970\,000} & \textbf{15,6\%} \\
\bottomrule
\end{tabular}
\end{table}

106 наблюдений --- нижняя граница, при которой градиентный бустинг хотя бы теоретически применим, но только при жёстком контроле сложности модели~\cite{makridakis2018m4}. Исходный набор включал 13 кандидатов-признаков; проведена двухэтапная отборка. На первом этапе по корреляционным и смысловым соображениям исключены: lag\_3, lag\_4 (слабая предсказательная сила при малой выборке), roll\_mean\_2 (нестабильный на малых выборках), month\_sin, month\_cos (дублируют week\_of\_year), is\_summer, is\_winter (избыточны при наличии week\_of\_year). На втором этапе оставшиеся признаки проверены пермутационным анализом~\cite{guyon2003introduction}: признак считается значимым, если его перестановка увеличивает WAPE более чем на 0,01. Итоговый набор из 6 признаков и результаты пермутационного анализа представлены в таблице~\ref{tab:feature_importance} и на рисунке~\ref{fig:feature_importance_chart}.

\begin{table}[H]
\centering
\caption{Результаты пермутационного анализа важности признаков}
\label{tab:feature_importance}
\begin{tabular}{lccc}
\toprule
\textbf{Признак} & \textbf{$\Delta$WAPE (перестановка)} & \textbf{Среднее} & \textbf{Std} \\
\midrule
share\_svetofor & $+0.235$ & 0.235 & 0.041 \\
roll\_mean\_4 & $+0.187$ & 0.187 & 0.029 \\
week\_of\_year & $+0.143$ & 0.143 & 0.038 \\
праздничные\_дни\_в\_месяце & $+0.094$ & 0.094 & 0.022 \\
lag\_2 & $+0.051$ & 0.051 & 0.018 \\
lag\_1 & $-0.012$ & $-0.012$ & 0.055 \\
\bottomrule
\end{tabular}
\end{table}

\begin{figure}[H]
\centering
\begin{tikzpicture}
\begin{axis}[
  xbar,
  width=11cm, height=5.5cm,
  bar width=0.42cm,
  xlabel={$\Delta$WAPE при перестановке},
  ytick=data,
  yticklabels={lag\_1, lag\_2, праздники, week\_of\_year, roll\_mean\_4, share\_svetofor},
  yticklabel style={font=\small},
  xmin=-0.06, xmax=0.28,
  xmajorgrids=true,
  grid style={dashed,gray!40},
  axis y line=left,
  axis x line=center,
  nodes near coords,
  nodes near coords align={horizontal},
  every node near coord/.append style={font=\scriptsize},
]
\addplot[fill=blue!55, draw=blue!70!black] coordinates {
  (-0.012,1) (0.051,2) (0.094,3) (0.143,4) (0.187,5) (0.235,6)
};
\end{axis}
\end{tikzpicture}
\caption{Пермутационная важность признаков (больше = важнее для модели)}
\label{fig:feature_importance_chart}
\end{figure}

Признак share\_svetofor (доля канала «Светофор» в продажах текущего месяца) является наиболее сильным предиктором --- это согласуется с доминирующей долей сети (85,3\%) в выручке. Отрицательная важность lag\_1 указывает на нестабильность оценки при малой выборке и чрезмерную подгонку под значение лага предыдущего месяца; признак оставлен, поскольку его значение близко к нулю.

\subsection{Методология кросс-валидации}

Оценка качества прогнозных моделей на временных рядах требует специальной схемы, исключающей утечку данных из будущего в прошлое~\cite{bergmeir2018note,kaufman2012leakage}. Применяется \textbf{rolling-origin} кросс-валидация (расширяющееся окно):
\begin{enumerate}
  \item Сначала обучение на первых $n_0$ наблюдениях, тест на следующих $h$.
  \item На каждой итерации обучающее окно расширяется на один шаг.
  \item Метрика усредняется по всем фолдам.
\end{enumerate}

Параметры схемы: $n_0 = 60$ (5 фолдов $\times$ 12 месяцев), $h = 3$ месяца (горизонт тестирования каждого фолда), число фолдов: 5. Схема показана на рисунке~\ref{fig:rolling_origin}. Дополнительно рассчитывается \textbf{recent WAPE} --- ошибка только на последнем фолде (самые свежие данные), позволяющая выявить деградацию качества при drift данных~\cite{widmer1996learning}.

\begin{figure}[H]
\centering
\begin{tikzpicture}[scale=0.82, every node/.style={font=\small}]
  \pgfmathsetmacro{\H}{0.52}
  \pgfmathsetmacro{\gap}{0.18}
  \pgfmathsetmacro{\testW}{1.8}
  \foreach \i in {1,...,5} {
    \pgfmathsetmacro{\y}{(\i-1)*(\H+\gap)}
    \pgfmathsetmacro{\trainW}{6 + (\i-1)*0.9}
    \fill[blue!30, draw=blue!60] (0,\y) rectangle (\trainW,\y+\H);
    \fill[orange!65, draw=orange!80] (\trainW,\y) rectangle (\trainW+\testW,\y+\H);
    \node[right] at (\trainW+\testW+0.15, \y+\H/2) {Фолд \i};
  }
  \node[blue!70, below] at (4, -0.25) {Обучение (расширяется)};
  \node[orange!80!black, below] at (9.3, -0.25) {Тест ($h = 3$ мес.)};
  \draw[-Stealth, thick] (-0.2,0) -- (-0.2, 5*(\H+\gap)+0.1) node[above]{};
  \draw[-Stealth, thick] (0,-0.6) -- (12.5,-0.6) node[right]{Время};
\end{tikzpicture}
\caption{Схема rolling-origin cross-validation (расширяющееся обучающее окно, фиксированный тестовый горизонт)}
\label{fig:rolling_origin}
\end{figure}

\subsection{Сравнительный анализ моделей и методологический артефакт}

Результаты всех протестированных моделей представлены в таблице~\ref{tab:demand_experiments} и на рисунке~\ref{fig:model_comparison_chart}.

\begin{table}[H]
\centering
\caption{Сравнение моделей прогноза спроса (CV WAPE, rolling-origin)}
\label{tab:demand_experiments}
\begin{tabular}{lccl}
\toprule
\textbf{Модель} & \textbf{CV WAPE} & \textbf{Std} & \textbf{Примечание} \\
\midrule
\textbf{moving\_average\_4} & \textbf{0.655} & 0.015 & \textbf{Победитель (бейзлайн)} \\
naive\_last & 0.712 & 0.028 & --- \\
seasonal\_naive & 0.731 & 0.041 & --- \\
moving\_average\_2 & 0.698 & 0.022 & --- \\
\midrule
HistGBM (sklearn) & 0.641 & 0.048 & Артефакт (см. ниже) \\
LightGBM (правильный) & 0.812 & 0.076 & Хуже бейзлайна \\
CatBoost & 0.783 & 0.064 & Хуже бейзлайна \\
Random Forest & 0.841 & 0.090 & Хуже бейзлайна \\
Extra Trees & 0.869 & 0.108 & Хуже бейзлайна \\
Ridge Regression & 0.723 & 0.039 & Хуже бейзлайна \\
ElasticNet & 0.741 & 0.043 & Хуже бейзлайна \\
\bottomrule
\end{tabular}
\end{table}

\begin{figure}[H]
\centering
\begin{tikzpicture}
\begin{axis}[
  xbar,
  width=12cm, height=8.5cm,
  bar width=0.38cm,
  xlabel={CV WAPE (меньше --- лучше)},
  ytick=data,
  yticklabels={
    Extra Trees,
    Random Forest,
    LightGBM,
    CatBoost,
    ElasticNet,
    Ridge,
    Seasonal Naive,
    Naive Last,
    MA(2),
    \textbf{MA(4)},
    HistGBM\textsuperscript{*}},
  yticklabel style={font=\small},
  xmin=0.58, xmax=0.93,
  xmajorgrids=true,
  grid style={dashed,gray!40},
  axis y line=left,
  axis x line=bottom,
  nodes near coords,
  nodes near coords align={horizontal},
  every node near coord/.append style={font=\tiny},
]
\addplot[fill=blue!45, draw=blue!65] coordinates {
  (0.869,1) (0.841,2) (0.812,3) (0.783,4)
  (0.741,5) (0.723,6) (0.731,7) (0.712,8)
  (0.698,9) (0.655,10) (0.641,11)
};
\draw[dashed, green!55!black, very thick]
  (axis cs:0.655,0.4) -- (axis cs:0.655,11.6);
\end{axis}
\end{tikzpicture}
\caption{CV WAPE всех протестированных моделей прогноза спроса. Пунктир --- результат победителя MA(4) = 0.655. * HistGBM --- методологический артефакт}
\label{fig:model_comparison_chart}
\end{figure}

Скользящее среднее с окном 4 (MA(4), WAPE = 0.655) побеждает все ML-модели. Это не означает, что «ML не работает» --- данных статистически недостаточно для надёжной оценки сложных моделей~\cite{makridakis2018m4}. При 106 наблюдениях MA(4) является оптимальным по критерию «сложность--качество».

В процессе экспериментов выявлен важный методологический артефакт~\cite{sculley2015hidden}: \texttt{HistGradientBoostingRegressor} (sklearn) показал WAPE = 0.641 --- лучше бейзлайна. При детальном разборе обнаружено, что sklearn-реализация использует \texttt{min\_samples\_leaf = 20} по умолчанию, тогда как LightGBM использует \texttt{min\_child\_samples = 5}. При 106 наблюдениях значение \texttt{min\_samples\_leaf = 20} фактически превращает бустинг в простое усреднение с неглубокими деревьями, имитируя скользящее среднее. При настройке LightGBM с аналогичным ограничением (\texttt{min\_child\_samples = 20}, \texttt{max\_depth = 3}) результаты двух библиотек сходятся; при \texttt{min\_child\_samples = 5} LightGBM переобучается --- WAPE = 0.812. Оптимальные параметры по CV: \texttt{n\_estimators = 50}, \texttt{max\_depth = 3}, \texttt{min\_child\_samples = 20}, что подтверждает необходимость жёсткой регуляризации при малых выборках.

Несмотря на итоговую победу бейзлайна, производственный прогноз рассчитывается по формуле планирования производства с учётом возвратов:

\begin{equation}
\text{план\_производства}_t = \frac{\hat{y}_t - \text{остаток}_t}{1 - r}
\end{equation}

где $\hat{y}_t$ --- прогнозируемый объём продаж, $\text{остаток}_t$ --- текущий складской остаток, $r$ --- ожидаемая доля возвратов. При $\hat{y}_t < \text{остаток}_t$ производство не требуется. Формула обеспечивает избегание как дефицита, так и перепроизводства.

\section{Прогнозирование возвратной продукции}

Возвраты в кондитерском производстве имеют двойственную природу: корректировка выручки и операционные потери (возвращённый товар, как правило, не может быть повторно реализован из-за истечения сроков годности или нарушения упаковки). Общий объём возвратов за период: \textbf{8\,300\,000 руб.} (около 39\,500 единиц продукции), 15,6\% от валовых продаж --- существенная статья потерь, поддающаяся прогнозированию.

Из данных 1С о корректировках реализации получены уровни возвратов в разрезе SKU (таблица~\ref{tab:return_rates} и рисунок~\ref{fig:return_rates_chart}).

\begin{table}[H]
\centering
\caption{Доля возвратов по SKU}
\label{tab:return_rates}
\begin{tabular}{lrr}
\toprule
\textbf{SKU} & \textbf{Return rate, \%} & \textbf{Интерпретация} \\
\midrule
Мадлены ассорти & 24,95\% & Критически высокий \\
Кекс с изюмом (Каприз) & 22,96\% & Очень высокий \\
ДИВНОЕ & 20,53\% & Высокий \\
Ажурная коса & 14,6\% & Умеренный \\
Песочное печенье & 5,2\% & Нормальный \\
\bottomrule
\end{tabular}
\end{table}

\begin{figure}[H]
\centering
\begin{tikzpicture}
\begin{axis}[
  xbar,
  width=12cm, height=5cm,
  bar width=0.45cm,
  xlabel={Доля возвратов, \%},
  ytick=data,
  yticklabels={Песочное печенье, Ажурная коса, ДИВНОЕ, Кекс с изюмом, Мадлены ассорти},
  yticklabel style={font=\small},
  xmin=0, xmax=30,
  xmajorgrids=true,
  grid style={dashed,gray!40},
  axis y line=left,
  axis x line=bottom,
  nodes near coords,
  nodes near coords align={horizontal},
  every node near coord/.append style={font=\small},
]
\addplot[fill=red!55, draw=red!70!black] coordinates {
  (5.2,1) (14.6,2) (20.53,3) (22.96,4) (24.95,5)
};
\draw[dashed, orange!80!black, very thick]
  (axis cs:12,0.4) -- (axis cs:12,5.55)
  node[above, font=\scriptsize, orange!80!black]{норма~12\%};
\end{axis}
\end{tikzpicture}
\caption{Доля возвратов по SKU, \%. Пунктир --- верхняя граница нормы для хлебобулочной продукции}
\label{fig:return_rates_chart}
\end{figure}

Мадлены ассорти и Кекс с изюмом демонстрируют патологически высокую долю возвратов (почти четверть отгруженного товара) --- нетипично даже для скоропортящейся продукции (средняя норма: 5--12\%). Вероятные причины: короткий срок годности при низкой оборачиваемости на полке «Светофора», завышенные объёмы поставки, проблемы с выкладкой.

Прогноз возвратов --- регрессионная задача с \textit{прерывистым} (intermittent) временным рядом~\cite{croston1972forecasting}: в большинстве дней возвраты равны нулю, редкие дни дают значительные суммы. Для таких рядов классические модели регрессии показывают систематическое завышение прогноза~\cite{syntetos2001bias}. Применён \textbf{CatBoost}~\cite{prokhorenkova2018catboost} с целевой переменной «объём возвратов в рублях» (дневная агрегация); признаки: ежедневный объём отгрузки, день недели, номер недели года, месяц, лаговые значения возвратов (1, 7, 14 дней), скользящее среднее за 7 дней. CatBoost выбран за нативную поддержку категориальных признаков (SKU, день недели) без ручного кодирования.

\begin{table}[H]
\centering
\caption{Сравнение моделей прогноза возвратов}
\label{tab:return_model}
\begin{tabular}{lccc}
\toprule
\textbf{Модель} & \textbf{CV WAPE} & \textbf{Recent WAPE} & \textbf{MAE, руб.} \\
\midrule
\textbf{CatBoost} & \textbf{0.709} & \textbf{0.668} & \textbf{тыс.} \\
naive\_last & 0.834 & 0.781 & --- \\
moving\_average\_7 & 0.812 & 0.743 & --- \\
Метод Кростона~\cite{croston1972forecasting} & 0.798 & 0.763 & --- \\
\bottomrule
\end{tabular}
\end{table}

CatBoost выигрывает у всех бейзлайнов: CV WAPE 0.709 против 0.798 у метода Кростона, разработанного специально для прерывистых рядов. Примечательно, что recent WAPE (0.668) лучше CV WAPE (0.709) --- улучшение качества на свежих данных является позитивным сигналом для производственного применения.

\section{Двухступенчатая hurdle-модель прогноза денежного потока}

Входящие и исходящие платежи имеют выраженную прерывистость: на большинстве дней нет ни одного крупного поступления или платежа определённого типа. Для таких распределений классические регрессионные модели систематически завышают прогноз~\cite{lambert1992zero}. Применяется \textbf{двухступенчатая (hurdle) модель}, разделяющая вероятность события и его интенсивность~\cite{lambert1992zero}:

\begin{equation}
\hat{y}_t = P(\text{non-zero}_t) \times \hat{y}^+_t
\end{equation}

\textbf{Ступень 1} (классификация): логистическая регрессия или градиентный бустинг предсказывает $P(\text{non-zero}_t)$ --- вероятность того, что в день $t$ будет хотя бы одна операция данного типа. \textbf{Ступень 2} (регрессия): модель предсказывает $\hat{y}^+_t$ --- ожидаемый объём платежей при условии их наличия. Признаковый набор включает: лаговые значения lag\_1, lag\_7, lag\_14, lag\_28; скользящие статистики roll\_mean\_7, roll\_mean\_14, roll\_std\_7; календарные переменные day\_of\_week, week\_of\_month, month, is\_month\_end, is\_month\_start; контекстные переменные (количество контрагентов за день, средняя сумма по контрагенту).

Сравнение моделей по входящим и исходящим платежам представлено в таблицах~\ref{tab:inflow_model} и~\ref{tab:outflow_model}.

\begin{table}[H]
\centering
\caption{Результаты прогноза входящих платежей (inflow)}
\label{tab:inflow_model}
\begin{tabular}{lccc}
\toprule
\textbf{Модель} & \textbf{CV WAPE} & \textbf{Recent WAPE} & \textbf{Примечание} \\
\midrule
\textbf{Hurdle LightGBM} & \textbf{0.237} & \textbf{0.122} & \textbf{Победитель} \\
Naive last & 0.429 & 0.371 & Бейзлайн \\
MA(7) & 0.398 & 0.342 & --- \\
Однофазная регрессия & 0.341 & 0.289 & Хуже hurdle \\
\bottomrule
\end{tabular}
\end{table}

\begin{table}[H]
\centering
\caption{Результаты прогноза исходящих платежей (outflow)}
\label{tab:outflow_model}
\begin{tabular}{lccc}
\toprule
\textbf{Модель} & \textbf{CV WAPE} & \textbf{Recent WAPE} & \textbf{Примечание} \\
\midrule
\textbf{Hurdle LightGBM} & \textbf{0.340} & \textbf{0.220} & \textbf{Победитель} \\
Naive last & 0.460 & 0.401 & Бейзлайн \\
MA(7) & 0.431 & 0.378 & --- \\
Однофазная регрессия & 0.389 & 0.311 & Хуже hurdle \\
\bottomrule
\end{tabular}
\end{table}

Hurdle-модель улучшила CV WAPE по inflow с 0.429 (naive last) до \textbf{0.237} --- снижение на 44,8\%. По outflow: с 0.460 до \textbf{0.340} --- снижение на 26,1\%. Recent WAPE по inflow (0.122) значительно лучше CV WAPE (0.237) --- позитивный тренд. Исходящие платежи прогнозируются хуже входящих (WAPE 0.340 vs 0.237): нерегулярные крупные расходы (аренда, квартальные налоги, разовые закупки) не имеют устойчивого временного паттерна. На основе обученных hurdle-моделей построен оперативный прогноз на ближайшие 30 дней: входящие платежи --- \textbf{1\,305\,991 руб.}, исходящие --- \textbf{1\,904\,507 руб.}, нетто-сальдо --- \textbf{$-598\,516$ руб.}. Отрицательный прогноз сальдо --- дополнительное свидетельство повторяемости кассовых разрывов; при наличии такого предупреждения за 30 дней у менеджмента достаточно времени для переговоров об отсрочке ряда платежей или ускорении поступлений.

\section{Детекция аномалий и мониторинг дрейфа моделей}

\subsection{Выявление аномальных транзакций}

Детектор аномалий решает задачу идентификации транзакций, значительно отличающихся от типичного паттерна по размеру, частоте или контрагенту~\cite{chandola2009anomaly}. В проекте применяются два комплементарных метода, обеспечивающих взаимную верификацию результатов.

\textbf{Метод 1: Robust z-score с MAD}~\cite{rousseeuw1993alternatives}. Классический z-score нечувствителен к распределениям с тяжёлыми хвостами, типичным для финансовых данных. Robust z-score использует медиану и медианное абсолютное отклонение (MAD) вместо среднего и стандартного отклонения:

\begin{equation}
z_\text{robust}(x) = \frac{x - \text{median}(X)}{c \cdot \text{MAD}(X)}, \quad c = 1.4826
\end{equation}

Константа $c = 1.4826$ обеспечивает совместимость оценки дисперсии с нормальным распределением~\cite{rousseeuw1993alternatives}. Операция считается аномальной при $|z_\text{robust}| > 3.5$.

\textbf{Метод 2: Isolation Forest}~\cite{liu2008isolation}. Алгоритм основан на свойстве аномалий изолироваться быстрее нормальных точек при случайном разбиении пространства признаков. Используются признаки: сумма операции, номер дня недели, номер дня месяца, тип контрагента, направление (\texttt{credit}/\texttt{debit}). Параметр \texttt{contamination = 0.02} (предполагаемая доля аномалий --- 2\%).

\begin{table}[H]
\centering
\caption{Результаты детекции аномалий}
\label{tab:anomaly_results}
\begin{tabular}{lr}
\toprule
\textbf{Показатель} & \textbf{Значение} \\
\midrule
Всего операций & 4\,086 \\
Флагов от robust z-score & 31 (0.76\%) \\
Флагов от Isolation Forest & 19 (0.46\%) \\
Согласованных флагов (оба метода) & 8 (0.20\%) \\
Аномальных дней (по дневному агрегату) & 3 \\
Итого операций с флагом аномалии & \textbf{42 (1.03\%)} \\
\bottomrule
\end{tabular}
\end{table}

Из 42 флагованных операций ручная проверка выявила два типа истинных аномалий: округлённые переводы на «ровные» суммы без реального контрагента (потенциальный вывод средств) и двойное поступление от клиента с последующим возвратом одной из сумм. Оба случая требуют разъяснения бухгалтерии.

\subsection{PSI-мониторинг и производственный контур}

Качество ML-моделей деградирует со временем по мере изменения распределения данных --- явление concept drift~\cite{widmer1996learning}. Для его мониторинга применяется \textbf{Population Stability Index (PSI)}~\cite{siddiqi2006credit}:

\begin{equation}
\text{PSI} = \sum_{i=1}^{B} \left( P_i^{\text{actual}} - P_i^{\text{expected}} \right) \times \ln\!\left( \frac{P_i^{\text{actual}}}{P_i^{\text{expected}}} \right)
\end{equation}

где $P_i^{\text{expected}}$ --- доля наблюдений в бакете $i$ на обучающем периоде, $P_i^{\text{actual}}$ --- на текущем. $B = 10$ бакетов. Пороги: PSI $< 0.10$ --- стабильно; $0.10 \le \text{PSI} < 0.25$ --- умеренный дрейф; PSI $\ge 0.25$ --- \textbf{переобучение необходимо}.

Расчёт PSI проведён по сравнению обучающего (май--октябрь 2025) и текущего (ноябрь 2025 -- апрель 2026) периодов (таблица~\ref{tab:psi_monitoring} и рисунок~\ref{fig:psi_chart}).

\begin{table}[H]
\centering
\caption{PSI мониторинг дрейфа моделей}
\label{tab:psi_monitoring}
\begin{tabular}{lccl}
\toprule
\textbf{Модель} & \textbf{PSI} & \textbf{Порог} & \textbf{Статус} \\
\midrule
Прогноз inflow & 1.032 & $\ge 0.25$ & \textbf{Переобучение} \\
Прогноз outflow & 1.165 & $\ge 0.25$ & \textbf{Переобучение} \\
Прогноз возвратов & 0.087 & $< 0.10$ & Стабильно \\
Прогноз спроса & 0.124 & $[0.10,0.25)$ & Мониторинг \\
\bottomrule
\end{tabular}
\end{table}

\begin{figure}[H]
\centering
\begin{tikzpicture}
\begin{axis}[
  xbar,
  width=10.5cm, height=5cm,
  bar width=0.48cm,
  xlabel={PSI},
  ytick=data,
  yticklabels={Возвраты, Спрос, Outflow, Inflow},
  yticklabel style={font=\small},
  xmin=0, xmax=1.4,
  xmajorgrids=true,
  grid style={dashed,gray!40},
  axis y line=left,
  axis x line=bottom,
  nodes near coords,
  nodes near coords align={horizontal},
  every node near coord/.append style={font=\small},
]
\addplot[fill=blue!50, draw=blue!70] coordinates {
  (0.087,1) (0.124,2) (1.165,3) (1.032,4)
};
\draw[dashed, orange!70!black, very thick]
  (axis cs:0.10,0.4) -- (axis cs:0.10,4.55)
  node[above, font=\scriptsize, orange!70!black, rotate=90, yshift=3pt]{0.10 мониторинг};
\draw[dashed, red!65, very thick]
  (axis cs:0.25,0.4) -- (axis cs:0.25,4.55)
  node[above, font=\scriptsize, red!65, rotate=90, yshift=3pt]{0.25 переобучение};
\end{axis}
\end{tikzpicture}
\caption{PSI мониторинг дрейфа. Значения inflow и outflow критически превышают порог переобучения}
\label{fig:psi_chart}
\end{figure}

PSI по inflow (1.032) и outflow (1.165) кратно превышают критический порог 0.25 --- прямое следствие сезонного паттерна: модели, обученные на весенне-летних данных, не описывают осенне-зимний период. PSI модели возвратов (0.087) остаётся стабильным: механизм формирования возвратов не зависит от сезона так же, как объёмы платежей. Для поддержания актуальности моделей реализован производственный контур~\cite{sculley2015hidden}: (1) ежедневный инкрементальный сбор данных через ETL-пайплайн; (2) автоматический расчёт PSI для каждой модели; (3) при PSI $\ge 0.25$ --- триггер переобучения; (4) переобучение на расширенном датасете; (5) валидация на holdout последнего месяца; (6) деплой при улучшении recent WAPE $\ge 5\%$; (7) автоматический управленческий отчёт с прогнозом, PSI и флагованными аномалиями.

\section{Ограничения и направления развития}

Текущее решение имеет четыре ключевых ограничения. \textbf{Малый размер выборки}: 106 наблюдений критически мало для надёжной оценки сложных моделей; каждый новый месяц добавляет 3 строки (3 активных SKU), и через 12--18 месяцев объём достигнет 140--160 строк --- минимального порога~\cite{makridakis2018m4}. \textbf{Единственный канал продаж}: 85\% выручки на «Светофоре» делает share\_svetofor главным предиктором, но одновременно ограничивает обобщаемость модели. \textbf{Частичная разметка}: 238 операций с флагом needs\_review требуют ручной проверки; до её завершения качество классификации не финальное~\cite{kaufman2012leakage}. \textbf{Нет складских остатков}: производственная формула предполагает знание остатков, которые в текущей версии вводятся вручную.

Приоритетные направления развития: интеграция складских данных из 1С для автоматизации производственной рекомендации; расширение признакового пространства за счёт внешних данных (индексы цен на сырьё, праздничный календарь, температурный фактор); накопление размеченного датасета транзакций для перехода от rule-based к ML-классификатору на собственных данных.
